% Generated by code/interviews-1/03-speed-drills/python/gen_drill.py. Do not edit. French copy by draft/pilot-quant-fr/ivgen_fr.py.
\ivfamily{Additionner et soustraire}{10, 11, 17, 24, 25, 28, 34, 36, 38, 40, 41, 42, 44, 49, 50, 53, 64, 66, 68, 83}{Travailler depuis la gauche, et ne pas écrire les retenues.}
\ivfamily{Tables de multiplication jusqu'à 19}{1, 2, 4, 5, 6, 7, 18, 20, 23, 26, 27, 31, 32, 33, 35, 37, 46, 48, 56, 58, 59, 63, 79, 80, 82}{Décomposer les nombres de 11 à 19~: $17 \times 8 = 80 + 56$.}
\ivfamily{Multiplier par un chiffre}{9, 12, 14, 15, 16, 54, 55, 57, 65, 70, 72, 84, 85, 86, 89}{Décomposer le plus grand nombre, les centaines d'abord.}
\ivfamily{Doubler et diviser par deux}{43, 52, 60, 62, 67, 69, 71, 77, 87, 90}{Procéder chiffre par chiffre depuis la gauche. Cela ne demande aucune réflexion.}
\ivfamily{Division exacte}{13, 30, 39, 45, 51, 73, 74, 75, 78, 88}{Penser la table de multiplication à l'envers, et décomposer un long nombre en parties faciles~: $476 \div 7 = 420 \div 7 + 56 \div 7$.}
\ivfamily{Pourcentages faciles}{3, 8, 19, 21, 22, 29, 47, 61, 76, 81}{$10\%$ déplace la virgule, $20\%$ double ce résultat, $25\%$ est un quart et $50\%$ une moitié.}
